%!TEX root = ../../main.tex

\chapter{Fazit und Ausblick}
\label{ch:fazit}

\section{Zusammenfassung der Ergebnisse}
\label{sec:zusammenfassung}

Ziel dieser Arbeit war eine experimentelle Evaluierung, inwieweit sich aktuelle \acp{LLM} für ausgewählte \ac{ETL}-Teilaufgaben eignen. Dazu wurde ein synthetischer Datensatz mit kontrolliert eingebauten Qualitätsproblemen erstellt, neun \ac{ETL}-Aufgaben in drei Kategorien und drei Schwierigkeitsstufen definiert und eine regelbasierte Vergleichspipeline implementiert, deren Ergebnis als Maßstab dient. Die Modelle wurden auf drei Untersuchungsebenen gemessen: in einer vollständig gekreuzten Faktormatrix aus Modell, Prompting-Strategie und Aufgabe, in einer verketteten Pipeline aus neun aufeinander aufbauenden Schritten sowie in einem sonst identischen isolierten Durchlauf, der jedem Schritt fehlerfreie Eingabedaten bereitstellt und damit eigene Leistung von geerbtem Fehler trennt. Ergänzend wurden zwei Ausführungsmodi gegenübergestellt: die Erzeugung von Verarbeitungscode und die unmittelbare Verarbeitung der Daten im Prompt. Alle Befunde wurden abschließend auf einem zweiten, während der Entwicklung unbekannten Datensatz mit mehreren Wiederholungen nachgeprüft.

Der über alle Ebenen hinweg stabilste Befund ist die Dominanz des Aufgabentyps. Zwischen dem besten und dem schwächsten kommerziellen Modell liegen rund zehn Prozentpunkte, zwischen der am besten und der am schlechtesten gelösten Aufgabe mehr als neunzig. Für ein konkretes \ac{ETL}-Vorhaben ist die Frage, welche Teilaufgaben einem Modell übertragen werden, damit erheblich folgenreicher als die Frage, welches Modell dafür gewählt wird -- ein Ergebnis, das der verbreiteten Orientierung an Modellranglisten eine begrenzte Aussagekraft bescheinigt.

Als zweiter Befund erwies sich der Ausführungsmodus als eigenständige und bislang unterschätzte Gestaltungsdimension. Der erwartete Vorteil der Modelle bei semantisch geprägten Aufgaben trat nur in der direkten Verarbeitung ein, in der das Modell sein Wissen auf jeden einzelnen Wert anwendet. Sobald dieselbe Aufgabe in generierten Code zu fassen war, fiel das Ergebnis auf das Niveau der regelbasierten Lösung zurück, weil das Modell sein Wissen in eine ebenfalls unvollständige Aufzählung oder einen numerischen Schwellenwert übersetzen musste. Umgekehrt scheiterte die direkte Verarbeitung an der abschließenden Aggregation über mehrere hundert Datensätze nahezu vollständig und kostete je verarbeiteter Zeile rund das Fünfzehnfache. Die beiden Modi sind damit nicht besser oder schlechter, sondern für unterschiedliche Aufgabenarten geeignet.

Der dritte Befund betrifft den Zuschnitt der Aufgabenstellung. Die zusammengesetzte Aufgabe aus Join, Bereinigung und Aggregation wurde im Code-Generierungs-Modus von keinem Modell und in keiner Strategie gelöst; dieselben Modelle lösten sie als Glied einer Kette mit fehlerfreier Eingabe vollständig korrekt. Nicht die Schwierigkeit der einzelnen Operation begrenzt die Leistungsfähigkeit, sondern die Zahl der Entscheidungen, die in einem Arbeitsschritt gleichzeitig zu treffen sind. Die Zerlegung einer Strecke ist damit der Eingriff mit dem größten Ertrag und zugleich der einzige, der ohne zusätzliche Kosten wirkt.

Der vierte Befund ergibt sich aus der Trennung von eigener Schrittleistung und Fehlerfortpflanzung. Der Endschritt der verketteten Strecke verlor je nach Modell zwischen 31 und 51~Prozentpunkte, die ausnahmslos aus vorgelagerten Schritten stammten und ihm selbst nicht zuzurechnen waren. Ein Schritt, der für sich betrachtet 94,6~Prozent erreichte und damit unauffällig wirkte, löste drei Schritte später einen vollständigen Abbruch aus. Für die Bewertung folgt daraus, dass Evaluierungen einzelner \ac{ETL}-Operationen die Leistungsfähigkeit in realen Strecken systematisch überschätzen, während rein verkettete Betrachtungen sie unterschätzen.

Der fünfte Befund korrigiert eine verbreitete Erwartung an die Prompt-Gestaltung, allerdings anders als zunächst angenommen. Die über alle Läufe gemittelte Rangfolge der Strategien erwies sich als nicht belastbar: Sie kehrte sich auf dem zweiten Datensatz um und ging, wie die Aufschlüsselung zeigte, nahezu vollständig auf das hardwarebegrenzte lokale Modell zurück, das unter den dreifach längeren Few-Shot-Prompts von 44,1 auf 24,0~Prozent einbrach. Unter den kommerziellen Modellen allein unterscheiden sich die Strategien über alle Aufgaben gemittelt kaum. Reproduzierbar ist stattdessen ein aufgabengebundenes Muster: Beispiele wirken bei Ähnlichkeitsentscheidungen, schrittweise Herleitung bei mehrstufigen Auswahlregeln, und bei sieben der neun Aufgaben wirkt keine der Strategien messbar. Zusätzlicher Kontext wirkt damit nicht durch seinen Umfang, sondern allein durch seine Passung zur Aufgabe -- und schadet, wo er über sie hinausweist.

Der sechste Befund ordnet diesem Ergebnis seinen eigentlichen Rang zu. Werden die neun Schritte zu drei Gruppen zusammengefasst und erhält jede Gruppe eine fehlerfreie Eingabe, so verschwindet der Strategieunterschied vollständig: Alle 27~Läufe lieferten ein verwertbares Ergebnis, die Genauigkeit lag zwischen 98,4 und 100~Prozent, die Spanne zwischen den Strategien bei einem halben Prozentpunkt. Was in der Einzelaufgaben-Matrix als Wirkung der Prompt-Gestaltung erscheint, ist demnach zu einem erheblichen Teil eine Wirkung des Aufgabenzuschnitts. Für die Praxis kehrt das die übliche Reihenfolge der Gestaltungsentscheidungen um: Zuerst ist die Aufgabe zuzuschneiden, erst danach der Prompt -- und ist der Zuschnitt richtig gewählt, verliert die Wahl der Strategie ihre Bedeutung weitgehend.

Der siebte Befund schließlich betrifft die Verlässlichkeit der Messung selbst. Im Bestätigungslauf schwankte der Endschritt derselben Konfiguration über fünf identische Wiederholungen zwischen 27,3 und 72,7~Prozent, während er im isolierten Durchlauf in jeder Wiederholung exakt denselben Wert lieferte. Die beobachtete Instabilität entstand folglich nicht aus der probabilistischen Textgenerierung selbst, sondern aus deren Verstärkung entlang der Kette -- und ist damit keine Eigenschaft des Modells, sondern eine Eigenschaft der Architektur, in die es eingebettet ist.

\section{Beantwortung der Forschungsfragen}
\label{sec:beantwortung-ff}

\paragraph{FF1.} Die erste Forschungsfrage richtete sich auf die Eignung aktueller \acp{LLM} für die Bearbeitung typischer \ac{ETL}-Teilaufgaben im Vergleich zu einer klassischen, regelbasierten Pipeline sowie auf die Bedingungen, unter denen je nach Aufgabentyp ein klassischer, \ac{KI}-gestützter oder hybrider Ansatz zu empfehlen ist.

Auf den ersten Teil der Frage lautet die Antwort, dass die Eignung nicht pauschal, sondern ausschließlich je Aufgabentyp und Ausführungsmodus zu beurteilen ist. Für die sechs formal beschreibbaren Aufgaben erreichen die kommerziellen Modelle mit gut gewähltem Prompt das Niveau der regelbasierten Pipeline, ohne dieses zu übertreffen; ihr Beitrag liegt dort nicht im Ergebnis, sondern darin, dass die Verarbeitungsvorschrift nicht zuvor formuliert werden musste. Für die semantisch geprägten Aufgaben übertreffen sie die regelbasierte Lösung, allerdings nur im Direkt-Modus und nur um wenige Prozentpunkte. Für die zusammengesetzte Aufgabe versagen alle untersuchten Ansätze einschließlich der regelbasierten Pipeline, solange sie ungeteilt gestellt wird. Die in H1 formulierte Erwartung ist damit in ihrem Kern bestätigt -- die starke Abhängigkeit vom Aufgabentyp, die Ebenbürtigkeit bei formalisierbaren Aufgaben, der Zuverlässigkeitsverlust bei zunehmender struktureller Komplexität und die Abwesenheit eines durchgängig überlegenen Ansatzes --, in einem Teilaspekt jedoch einzuschränken: Der erwartete Vorteil bei semantischen Aufgaben tritt nicht generell ein, sondern ist an den Ausführungsmodus gebunden und bleibt im primär untersuchten Modus der Code-Generierung aus.

Auf den zweiten Teil der Frage folgt die Empfehlung eines hybriden Vorgehens, das in Abschnitt~\ref{sec:hybride-architekturen} ausgeführt ist: regelbasierte Ausführung für die syntaktisch beschreibbaren Schritte, ein Modell im Direkt-Modus für die wenigen semantisch geprägten Schritte, regelbasierte oder codegenerierte Ausführung für die abschließende Aggregation, und dies durchgehend in kleinteilig zugeschnittenen Schritten mit einer inhaltlichen Prüfung nach jedem einzelnen. Für die untersuchte Strecke bedeutet dies, dass sieben der neun Schritte regelbasiert auszuführen wären und lediglich zwei ein Modell benötigen. Ein durchgängiger \ac{LLM}-Einsatz ist auf Grundlage dieser Ergebnisse nicht zu empfehlen, ein durchgängiger Verzicht ebenso wenig.

\paragraph{FF2.} Die zweite Forschungsfrage richtete sich auf den Einfluss der Prompting-Strategie auf die Ergebnisqualität und auf die daraus abzuleitende Empfehlung.

Der Einfluss ist erheblich, aber ungleich verteilt und stark vom Aufgabenzuschnitt abhängig. Bei sieben der neun Aufgaben war zwischen den Strategien kein Unterschied messbar; bei den beiden anspruchsvollen Entdoppelungsaufgaben entschied die Wahl über bis zu 55~Prozentpunkte, und die schema-angereicherte Variante senkte die Genauigkeit dort auf ein Viertel. Wo die Prompt-Gestaltung überhaupt wirkt, ist ihr Gewicht demnach mit dem der Modellwahl vergleichbar.

H2 ist in ihrer allgemeinen Form zu verwerfen, allerdings mit einer Begründung, die von der zunächst naheliegenden abweicht. Aufwändigere Strategien verbessern die Ergebnisqualität gegenüber dem Zero-Shot-Prompt nicht durchgängig -- aber sie verschlechtern sie unter den kommerziellen Modellen auch nicht durchgängig. Die über alle Aufgaben gemittelte Rangfolge erwies sich schlicht als nicht belastbar: Sie kehrte sich zwischen den beiden Datensätzen um und ging fast vollständig auf das hardwarebegrenzte lokale Modell zurück. Zutreffend bleibt der erwartete abnehmende Grenznutzen in verschärfter Form -- bei unpassendem Zuschnitt ist er negativ --, ebenso die erwartete Aufgabenabhängigkeit, die sich über beide Datensätze hinweg reproduzierte.

Die Empfehlung, die daraus folgt, betrifft weniger die Wahl der Strategie als die Reihenfolge der Entscheidungen. Wird dieselbe Verarbeitung in angemessen zugeschnittenen Gruppen statt als überladene Einzelaufgabe gestellt, liegen alle drei Strategien zwischen 99,0 und 99,5~Prozent, und die Frage nach der besten Strategie erübrigt sich. Zuerst ist daher die Aufgabe zu zerlegen; bleibt danach eine Aufgabe übrig, die sich nicht weiter teilen lässt, gelten die beiden reproduzierten Zuordnungen: Beispiele bei Ähnlichkeitsentscheidungen, schrittweise Herleitung bei mehrstufigen Auswahlregeln, im Übrigen der einfachste und billigste Prompt. Zusätzlicher Kontext ist ausschließlich als Reaktion auf einen konkret beobachteten Mangel zu ergänzen -- so behob die Angabe der Eingabe-Datentypen ein stilles Versagen der Spaltenauswahl --, niemals vorsorglich.

\paragraph{Reproduzierbarkeit.} Die querschnittlich erhobene Reproduzierbarkeitsdimension fügt beiden Antworten eine Einschränkung hinzu. Die niedrig gewählte Temperatur begrenzte die Variabilität wirksam, beseitigte sie aber nicht: In gut einem Viertel der Konfigurationen führte derselbe Prompt bei demselben Modell zu messbar unterschiedlichen Ergebnissen, und die Instabilität konzentrierte sich auf genau jene Aufgaben, die auch inhaltlich am schwersten fielen. Gerade dort, wo ein Ergebnis besonders sorgfältig zu prüfen wäre, ist es zugleich am wenigsten verlässlich wiederholbar. Da die Verstärkung entlang der Kette sich als der dominierende Beitrag zu dieser Streuung erwies, ist sie durch Festschreibung des geprüften Codes und der frühen Zwischenergebnisse jedoch weitgehend beherrschbar.

\section{Limitationen der Arbeit}
\label{sec:limitationen}

Die vorstehenden Aussagen unterliegen mehreren Einschränkungen, die bei ihrer Übertragung zu berücksichtigen sind.

Die erste und gewichtigste betrifft den Datensatz. Er ist synthetisch erzeugt, und seine Qualitätsprobleme entstammen einer bekannten, endlichen Menge von Mustern. Diese Anlage war notwendig, um eine eindeutige Soll-Lösung zu besitzen und die Ansätze überhaupt vergleichbar messen zu können, begünstigt jedoch systematisch die regelbasierte Vergleichspipeline: Deren Zuordnungstabelle und Normalisierungsregel wurden gegen genau diese Muster entworfen und erreichen deshalb Werte, die in einem realen Bestand mit unbekannten Abweichungsmustern nicht zu erwarten wären. Der in dieser Arbeit gemessene Vorsprung der regelbasierten Lösung ist damit als Obergrenze zu lesen, und das Verhältnis verschiebt sich in realen Beständen zugunsten der semantischen Verfahren -- um welchen Betrag, lässt sich aus diesen Daten nicht bestimmen.

Die zweite Einschränkung betrifft die Größenordnung. Der Datensatz umfasst einige hundert Kunden und einige tausend Bestellungen, was für die inhaltliche Bewertung ausreicht, aber keine Aussage über das Verhalten bei produktionstypischen Datenmengen erlaubt. Für den Code-Generierungs-Modus ist dies unkritisch, da die Datenmenge dort nicht durch das Modell hindurchgeht; für den Direkt-Modus hingegen ist es die bestimmende Grenze, die im Rahmen dieser Arbeit bereits bei wenigen hundert Datensätzen eine Halbierung des Datensatzes erzwang.

Eng damit verbunden ist eine Einschränkung zum lokal betriebenen Modell. Sein Rückstand ist maßgeblich durch den verfügbaren Grafikspeicher bestimmt, der die Größe des ausführbaren Modells begrenzt; gemessen wurde insofern nicht allein Modellqualität, sondern auch die Ausstattung des ausführenden Rechners. Aus diesem Grund wurde es aus den Vergleichen der Prompting-Strategien, der Pipeline und des Bestätigungslaufs herausgenommen -- eine Entscheidung, die die Vergleichbarkeit dieser Ebenen sichert, zugleich aber bedeutet, dass für lokal betreibbare Modelle nur die Ergebnisse der Faktormatrix vorliegen. Ob ein größeres, auf leistungsfähigerer Hardware ausführbares Modell den Abstand schließt, bleibt hier offen; die Fehlerverteilung -- weit überwiegend nicht lauffähiger Code statt inhaltlicher Fehler -- legt es nahe, belegt es aber nicht.

Die dritte Einschränkung betrifft die Momentaufnahme. Untersucht wurden vier Modelle zu einem bestimmten Zeitpunkt in einem Feld, das sich in Monatsabständen verändert. Die exakte Versionskennung wurde bei jeder Ausführung protokolliert, sodass die Läufe dokumentiert bleiben; da die kommerziellen Modelle jedoch serverseitig gepflegt werden, lässt sich ihr Verhalten weder einfrieren noch nach einer Modellaktualisierung wiederherstellen. Die konkreten Zahlenwerte sind daher nicht als dauerhaft gültig zu betrachten -- die Wirkungszusammenhänge, aus denen sie entstehen, hingegen sehr wohl, da sie in ihren Ursachen benannt und im Bestätigungslauf repliziert wurden.

Die vierte Einschränkung betrifft die Zahl der Wiederholungen. Für die Faktormatrix wurden drei Wiederholungen je Konfiguration ausgeführt, für die verkettete Pipeline fünf. Angesichts der beobachteten Streuung -- allein der Endschritt schwankte über fünf Wiederholungen um bis zu 45~Prozentpunkte -- ist dies eine schmale Grundlage für Aussagen über einzelne Konfigurationen. Die Befunde dieser Arbeit sind deshalb bewusst auf Ebene der Muster formuliert, die über Modelle und Aufgaben hinweg wiederkehren, und nicht auf Ebene einzelner Zahlenwerte.

Die fünfte Einschränkung betrifft den Zuschnitt der Untersuchung. Betrachtet wurde die Transformationsphase des \ac{ETL}-Prozesses; Extraktion und Ladevorgang blieben ebenso außen vor wie Aspekte des produktiven Betriebs, etwa Fehlerbehandlung, Überwachung, inkrementelle Verarbeitung oder Datenschutzanforderungen beim Übertragen von Daten an externe Anbieter. Gerade der letzte Punkt ist für die praktische Anwendbarkeit erheblich und wurde hier nur insoweit berührt, als mit dem lokal betriebenen Modell eine Option ohne Datenabfluss mituntersucht wurde.

Schließlich ist die Rolle der prüfenden Instanz zu benennen. Die Untersuchung hat die Modelle weitgehend autonom arbeiten lassen und ihre Ergebnisse gegen eine bekannte Soll-Lösung bewertet. Der in der Praxis dominierende Einsatz als Assistenz einer Fachkraft, die den erzeugten Code liest, korrigiert und verantwortet, wurde damit nicht abgebildet. Die in dieser Arbeit gemessenen Werte beschreiben folglich, was ohne menschliche Nachbearbeitung entsteht -- was zugleich die Untergrenze dessen markiert, was im assistierten Einsatz zu erwarten ist.

\section{Ausblick auf weitere Forschung}
\label{sec:ausblick}

Aus den Ergebnissen und ihren Grenzen ergeben sich mehrere Anschlussfragen.

Am unmittelbarsten anschlussfähig ist die Übertragung auf reale Datenbestände. Da der Vorsprung der regelbasierten Pipeline maßgeblich daraus resultiert, dass ihre Regeln gegen bekannte Abweichungsmuster entworfen wurden, wäre eine Wiederholung des Vergleichs auf einem gewachsenen Bestand mit unbekannten Mustern der aussagekräftigste nächste Schritt. Zu erwarten ist, dass sich das Verhältnis zugunsten der semantischen Verfahren verschiebt; die Größe dieser Verschiebung ist die eigentlich interessante Messgröße und zugleich diejenige, die über die wirtschaftliche Tragfähigkeit eines \ac{LLM}-Einsatzes entscheidet.

Ein zweiter Ansatzpunkt ist die in Abschnitt~\ref{sec:hybride-architekturen} skizzierte hybride Architektur, die aus den Ergebnissen abgeleitet, aber nicht mehr als Ganzes implementiert und gemessen wurde. Eine Umsetzung, die regelbasierte Schritte, gezielt eingesetzte Direkt-Modus-Schritte und eine Prüfung nach jedem Schritt tatsächlich verbindet, ließe sich unmittelbar gegen die hier gemessenen rein regelbasierten und rein \ac{KI}-gestützten Varianten stellen und würde zeigen, ob die erwartete Kombination der Stärken auch in der Ausführung eintritt.

Der dritte Ansatzpunkt betrifft die Fehlerfortpflanzung, deren Messung diese Arbeit auf den Vergleich zweier Durchläufe stützt. Offen bleibt, welche Eigenschaften eines Zwischenergebnisses seine Verwendbarkeit in nachgelagerten Schritten bestimmen. Die Ergebnisse legen nahe, dass nicht der Umfang, sondern die Art der verbliebenen Abweichung entscheidet -- 145 unveränderte Zeitstempel unter 540~Datumsangaben brachten die Kette zum Abbruch, während gleichmäßig verteilte kleinere Ungenauigkeiten folgenlos blieben. Eine Kennzahl, die diese Unterscheidung abbildet, wäre für die Überwachung \ac{KI}-gestützter Strecken erheblich wertvoller als die übliche Genauigkeit je Schritt.

Ein vierter Ansatzpunkt ergibt sich aus dem Befund zum Ausführungsmodus. Da die direkte Verarbeitung bei semantischen Einzelentscheidungen überlegen, bei Aggregationen über viele Datensätze jedoch unbrauchbar ist, liegt eine Kombination nahe, in der das Modell die semantische Entscheidung je Wert trifft und deren Ergebnis anschließend deterministisch weiterverarbeitet wird -- etwa als generierte Zuordnungstabelle statt als generierte Aufzählungslogik. Ob eine solche Zwischenform den Genauigkeitsvorteil des Direkt-Modus bei dessen Kostenstruktur vermeidet, ist eine unmittelbar prüfbare Frage.

Schließlich verdient der in Abschnitt~\ref{sec:limitationen} angesprochene assistierte Einsatz eigene Aufmerksamkeit. Diese Arbeit hat gemessen, was Modelle ohne menschliche Nachbearbeitung leisten; die praktisch relevantere Frage ist, wie viel Aufwand eine Fachkraft benötigt, um ein erzeugtes Ergebnis auf ein verantwortbares Niveau zu bringen, und ob dieser Aufwand geringer ausfällt als der einer eigenständigen Implementierung. Gerade der in dieser Arbeit dokumentierte Fall des stillen Versagens -- eines Codes, der fehlerfrei durchläuft, schemakonform ausgibt und dennoch wirkungslos bleibt -- legt nahe, dass diese Frage nicht allein über die Ergebnisqualität, sondern wesentlich über den Prüfaufwand zu beantworten ist.
